\documentclass[12pt]{article}
\usepackage[utf8]{inputenc}
\usepackage[french]{babel}
\usepackage[style=numeric, sorting=none]{biblatex}

\addbibresource{computer science unplugged.bib}

\nocite{alayrangues_informatique_2017}
\nocite{bell_computer_2012}
\nocite{nishida_new_2008}
\nocite{sendurur_investigation_2019}
\nocite{sharif_university_of_technology_computational_2017}
\nocite{srivastava_assessing_2019}

\title{A Computer Science Unplugged Activity: CityMap}
\author{Merve Yildiz\textsuperscript{1} and Hasan Kara\textsuperscript{2}}
\date{Karadeniz Technical University, Turkey\textsuperscript{1} \\Trabzon University, Turkey\textsuperscript{2}}

\begin{document}
\maketitle

\section{Introduction}
In recent years, the concept of computational thinking has come to the fore as computer science has become a multidisciplinary 
field and the teaching of programming has spread rapidly at the K-12 education. Computational thinking, which is described as 
one of the 21st century skills \cite{sharif_university_of_technology_computational_2017} , is seen as a basic literacy skill for 
digital age learners \cite{nishida_new_2008}. When the literature is examined, the concept of computational thinking was first 
put forward by \cite{bell_computer_2012} as computational ideas. Although there are different definitions of the concept of 
computational thinking, it is emphasized in the literature that computational thinking is mostly related to problem solving and 
algorithmic thinking \cite{sendurur_investigation_2019, srivastava_assessing_2019}. In this context, for this study, 
computational thinking is considered as writing algorithms using step by step instructions towards the solution of a problem 
\cite{alayrangues_informatique_2017}.


\printbibliography[title={References}]

\end{document}
